% ---- pg-553 （书页 541） ----

\begin{mvfigure}
\mvinc{pg-553}{0.121,0.117,0.878,0.795}
\phantomsection\label{fig:12.11}\mvcap{图 12.11\quad 109 种纯麦芽苏格兰威士忌之间相似性的树图}
\end{mvfigure}

\vspace{0.4em}
\noindent\rule{\linewidth}{0.4pt}

{\zihao{6}
\noindent * 树图中标有字母 A 至 L 的 12 个组是经研究者识别后认为是相似的. 随后的分析表明, 这 12 个组受到某个很大的地理成分的支配.
也就是说, 位置互相靠近的那些酒厂倾向于生产具有类似特性的威士忌. 于是这些研究人员得出结论: “与地理特征的关系已被证实.
分析结果支持如下假设: 威士忌不仅受到厂家生产秘方和传统的影响, 而且还受到与地区有关的诸如\par}
